% !TEX root = ../main.tex
\begin{ouverture}[lunette]{24}{Lunette astronomique}{Partie 4 · Ondes et signaux}
{Lentille convergente\enspace·\enspace Lunette afocale\enspace·\enspace Faisceau émergent\enspace·\enspace Grossissement}
\objectif{Représenter le schéma d'une lunette afocale modélisée par deux lentilles minces convergentes ; identifier
l'objectif et l'oculaire.}
\objectif{Représenter le faisceau émergent issu d'un point objet situé à l'infini et traversant une lunette afocale.}
\objectif{Établir l'expression du grossissement d'une lunette afocale.}
\objectif{Exploiter les données caractéristiques d'une lunette commerciale.}
\objectif{Calculer en Python le diamètre apparent d'un astre observé avec différents oculaires.}
\prerequis{Lentille mince convergente : foyers, distance focale, vergence, construction d'images (classe de première).}
\prerequis{Relation de conjugaison et grandissement (classe de première).}
\prerequis{Tangente d'un angle, approximation des petits angles (mathématiques).}
\end{ouverture}
\essentiel
\essSection{1}{Image d'un objet à l'infini}
\begin{propriete}[Lentille convergente et objet à l'infini]
Les rayons issus d'un point très éloigné, comme une étoile, arrivent sur une lentille parallèles entre eux. Après une
lentille convergente de distance focale $f'$, ils convergent en un point du \textbf{plan focal image}, situé sur le rayon
qui passe sans déviation par le centre optique $O$. Si le faisceau est incliné d'un angle $\theta$ sur l'axe optique, ce
point est à la distance $f'\tan\theta$ de l'axe. Réciproquement, les rayons issus d'un point du plan focal objet ressortent
parallèles entre eux.
\end{propriete}
\schema{24-lentille}
\essSection{2}{La lunette afocale}
\begin{definition}[Lunette astronomique]
Une lunette astronomique est modélisée par deux lentilles minces convergentes de même axe optique :
\begin{itemize}
\item l'\textbf{objectif} $L_1$, du côté de l'objet, de grande distance focale $f'_1$ ;
\item l'\textbf{oculaire} $L_2$, du côté de l'œil, de petite distance focale $f'_2$.
\end{itemize}
La lunette est \textbf{afocale} lorsque le foyer image de l'objectif est confondu avec le foyer objet de l'oculaire :
$F'_1=F_2$. Elle donne alors d'un objet à l'infini une image à l'infini, que l'œil observe sans accommoder, donc sans
fatigue. Les deux lentilles sont distantes de $f'_1+f'_2$.
\end{definition}
\begin{propriete}[Marche des rayons]
L'objectif donne d'un astre une \textbf{image intermédiaire} $A_1B_1$, réelle et renversée, dans son plan focal image.
Cette image, située dans le plan focal objet de l'oculaire, joue pour lui le rôle d'objet : l'oculaire en donne une
image à l'infini. Le faisceau émergent est formé de rayons parallèles au rayon qui joint $B_1$ au centre optique $O_2$ de
l'oculaire.
\end{propriete}
\schema{24-lunette}
\essSection{3}{Le grossissement}
\begin{definition}[Grossissement]
Le \textbf{grossissement} d'un instrument est le rapport $G=\dfrac{\theta'}{\theta}$, où $\theta$ est l'angle sous lequel on voit
l'objet à l'œil nu et $\theta'$ l'angle sous lequel on voit son image à travers l'instrument.
\end{definition}
\begin{savoirfaire}[Établir l'expression du grossissement]
Dans le triangle $O_1F'_1B_1$ : $\tan\theta=\dfrac{A_1B_1}{f'_1}$ ; dans le triangle $O_2F_2B_1$ : $\tan\theta'=\dfrac{A_1B_1}{f'_2}$. Pour de
petits angles, $\tan\theta\approx\theta$ et $\tan\theta'\approx\theta'$, d'où
\[G=\frac{\theta'}{\theta}=\frac{f'_1}{f'_2}.\]
Le grossissement est d'autant plus grand que l'objectif a une grande distance focale et l'oculaire une petite.
\end{savoirfaire}
\begin{exemple}
Une lunette commerciale « 70/700 » a un objectif de \qty{70}{mm} de diamètre et de \qty{700}{mm} de distance focale. Avec un
oculaire de \qty{25}{mm}, $G=28$ ; avec un oculaire de \qty{10}{mm}, $G=70$. La Lune, vue à l'œil nu sous
$\theta=0{,}52^\circ=\qty{9,1e-3}{rad}$, donne une image intermédiaire de $f'_1\tan\theta=\qty{6,4}{mm}$. Le grossissement maximal
utile d'une lunette est de l'ordre de deux fois le diamètre de son objectif exprimé en millimètres, ici 140.
\end{exemple}
\begin{codecote}[0.65]
\begin{pythoncol}
from math import atan, tan, degrees, radians
f1 = 700                      # mm, objectif
theta = radians(0.52)         # diamètre apparent
for f2 in [25, 10, 6]:        # mm, oculaires
    G = f1 / f2
    theta2 = atan(G * tan(theta))
    print(f2, round(G), round(degrees(theta2), 1))
\end{pythoncol}
\begin{consolecol}
25 28 14.3
10 70 32.4
6 117 46.6
\end{consolecol}
\cote
\fontsize{9.5}{12.5}\selectfont
\capacite{Capacité numérique}\par\smallskip
Pour chaque oculaire, le programme calcule le grossissement et l'angle, en degrés, sous lequel la Lune apparaît dans la
lunette, à partir de $\tan\theta'=\dfrac{f'_1}{f'_2}\tan\theta$. Avec l'oculaire de \qty{6}{mm}, la relation $\theta'=G\,\theta$ donnerait
$60{,}7^\circ$ : l'approximation des petits angles n'est plus valable.
\end{codecote}
\begin{attention}
Le grossissement compare des angles ; il ne faut pas le confondre avec le grandissement, rapport de deux tailles. La
lunette donne une image renversée, ce qui est sans importance en astronomie. Elle ne grossit pas les étoiles, trop
lointaines pour apparaître autrement que comme des points, mais elle collecte d'autant plus de lumière que son objectif
est grand.
\end{attention}
\begin{experience}[Réaliser une maquette de lunette]
Sur un banc d'optique, on place un objectif de vergence \num{+2}~$\delta$ et un oculaire de vergence \num{+8}~$\delta$, distants de
\qty{62,5}{cm}. Un objet placé au foyer d'une lentille auxiliaire joue le rôle d'objet à l'infini. On visualise l'image
intermédiaire sur un écran placé dans le plan focal image de l'objectif, puis on retire l'écran pour observer à travers
l'oculaire. En mesurant la taille de l'image intermédiaire et les angles, on vérifie que $G=\dfrac{f'_1}{f'_2}=4$.
\end{experience}
\begin{rediger}
Pour construire le faisceau émergent, tracer d'abord le rayon incident qui passe par $O_1$ : il donne $B_1$ dans le plan
focal $F'_1=F_2$. Tous les rayons du faisceau convergent en $B_1$. Tracer ensuite le rayon $B_1O_2$ : les rayons émergents lui
sont parallèles.
\end{rediger}
\exercices

\rubrique{Découverte}
\begin{blocLivre}{ex-24-1}
\begin{exercice}[1]{Le vocabulaire de la lunette}
Une lunette astronomique comporte deux lentilles convergentes de distances focales \qty{900}{mm} et \qty{20}{mm}.
\begin{questions}
\item Laquelle constitue l'objectif ? l'oculaire ? Justifier.
\item Que signifie « lunette afocale » ? Quel est l'intérêt de ce réglage pour l'observateur ?
\item Quelle est alors la distance entre les deux lentilles ?
\item Calculer le grossissement de la lunette.
\end{questions}
\end{exercice}
\end{blocLivre}
\begin{blocLivre}{ex-24-2}
\begin{exercice}[1]{L'image de la Lune}
La Lune est vue depuis la Terre sous un angle de $0{,}52^\circ$. On en forme l'image avec une lentille convergente.
\begin{questions}
\item Où se forme l'image de la Lune ? Est-elle réelle ou virtuelle, droite ou renversée ?
\item Calculer le diamètre de cette image pour une lentille de distance focale \qty{1,0}{m}.
\item Même question pour le téléobjectif d'un appareil photo, de distance focale \qty{300}{mm}, puis pour une lunette de
\qty{2000}{mm} de distance focale.
\item Pourquoi les astronomes amateurs recherchent-ils des instruments de grande distance focale ?
\end{questions}
\end{exercice}
\end{blocLivre}
\begin{blocLivre}{ex-24-3}
\begin{exercice}[1]{Changer d'oculaire}
Une lunette a un objectif de \qty{1000}{mm} de distance focale. Elle est livrée avec trois oculaires de \qty{40}{mm},
\qty{20}{mm} et \qty{9}{mm}.
\begin{questions}
\item Rappeler l'expression du grossissement d'une lunette afocale.
\item Calculer le grossissement obtenu avec chaque oculaire.
\item Quel oculaire choisir pour observer la Lune entière ? les détails d'un cratère ?
\item Faut-il déplacer l'oculaire quand on en change ? Justifier.
\end{questions}
\end{exercice}
\end{blocLivre}
\begin{blocLivre}{ex-24-4}
\begin{exercice}[1]{Construire le faisceau émergent}
Le schéma représente une lunette afocale et un faisceau de rayons parallèles issu d'un point situé à l'infini.
\schema{24-ex04}
\begin{questions}
\item Placer l'image intermédiaire $B_1$ du point à l'infini. Justifier la construction.
\item Tracer la suite des trois rayons jusqu'à l'oculaire.
\item Construire le faisceau émergent.
\item Mesurer sur le schéma les angles $\theta$ et $\theta'$, puis vérifier que leur rapport est égal à $f'_1/f'_2$.
\end{questions}
\end{exercice}
\end{blocLivre}
\begin{blocLivre}{ex-24-5}
\begin{exercice}[1]{Vergences et distances focales}
Une petite lunette comporte un objectif de vergence \num{+1,25}~$\delta$ et un oculaire de vergence \num{+40}~$\delta$.
\begin{questions}
\item Rappeler la relation entre vergence et distance focale. Calculer les distances focales des deux lentilles.
\item Calculer le grossissement de la lunette.
\item Calculer la longueur de la lunette réglée en position afocale.
\item Exprimer le grossissement à l'aide des vergences.
\end{questions}
\end{exercice}
\end{blocLivre}
\rubrique{Entraînement}
\begin{blocLivre}{ex-24-6}
\begin{exercice}[2]{La Lune dans la lunette}
On observe la Lune, de diamètre apparent $0{,}52^\circ$, avec une lunette 70/700 munie d'un oculaire de \qty{10}{mm}.
\begin{questions}
\item Calculer le grossissement.
\item Calculer, avec l'approximation des petits angles, l'angle sous lequel la Lune apparaît dans la lunette. Calculer
l'angle exact à l'aide des tangentes. Commenter.
\item Calculer le diamètre de l'image intermédiaire de la Lune.
\item Le champ de l'oculaire, c'est-à-dire l'angle maximal sous lequel l'œil voit à travers lui, vaut $50^\circ$. Voit-on la
Lune entière ?
\end{questions}
\end{exercice}
\end{blocLivre}
\begin{blocLivre}{ex-24-7}
\begin{exercice}[2]{Voir les cratères lunaires}
L'œil distingue deux points vus sous un angle d'au moins \qty{3e-4}{rad}. On observe la Lune, à \qty{3,84e5}{km}, avec la
lunette 70/700 et un oculaire donnant $G=70$.
\begin{questions}
\item Quelle est la taille du plus petit cratère visible à l'œil nu ?
\item Quel angle minimal l'œil doit-il voir à travers la lunette ? Quelle est la taille du plus petit cratère visible ?
\item La diffraction par l'objectif limite la résolution à $\theta_\mathrm{min}=1{,}22\,\dfrac{\lambda}{D}$, avec $\lambda=\qty{550}{nm}$ et
$D$ le diamètre de l'objectif. Calculer la taille correspondante sur la Lune.
\item Expliquer pourquoi il est inutile d'augmenter indéfiniment le grossissement.
\end{questions}
\end{exercice}
\end{blocLivre}
\begin{blocLivre}{ex-24-8}
\begin{exercice}[2]{Programmer le grossissement}
On reprend le programme de l'essentiel.
\begin{questions}
\item \capacite{Capacité numérique}Pourquoi le programme convertit-il l'angle en radians ? Que calcule la ligne
\code{theta2 = atan(G * tan(theta))} ?
\item Justifier cette ligne à l'aide de l'image intermédiaire.
\item Modifier le programme pour qu'il affiche aussi la valeur approchée $G\,\theta$, en degrés. Jusqu'à quel grossissement
l'écart reste-t-il inférieur à 5~\% ?
\item On observe Jupiter, de diamètre apparent $45''$ (secondes d'arc). Quelle ligne faut-il modifier ?
\end{questions}
\end{exercice}
\end{blocLivre}
\begin{blocLivre}{ex-24-9}
\begin{exercice}[2]{Retrouver les distances focales}
Une lunette afocale mesure \qty{1,05}{m} entre l'objectif et l'oculaire et grossit 20 fois.
\begin{questions}
\item Écrire les deux relations vérifiées par $f'_1$ et $f'_2$.
\item En déduire les distances focales de l'objectif et de l'oculaire.
\item Calculer leurs vergences.
\item On remplace l'oculaire par un autre de \qty{2,5}{cm}. Quel est le nouveau grossissement ? Que faut-il régler ?
\end{questions}
\end{exercice}
\end{blocLivre}
\begin{blocLivre}{ex-24-10}
\begin{exercice}[2]{Regarder par le mauvais bout}
On regarde à travers la lunette de l'exercice~1 ($f'_1=\qty{900}{mm}$, $f'_2=\qty{20}{mm}$) en plaçant l'œil du côté de l'objectif.
\begin{questions}
\item Quelle lentille joue maintenant le rôle d'objectif ? d'oculaire ?
\item Le réglage reste-t-il afocal ?
\item Calculer le nouveau grossissement. Comment voit-on les objets ?
\item Citer un dispositif courant qui utilise ce principe.
\end{questions}
\end{exercice}
\end{blocLivre}
\begin{blocLivre}{ex-24-11}
\begin{exercice}[2]{Les jumelles 10 × 50}
Des jumelles portent l'indication « 10 × 50 » et « champ : $6{,}5^\circ$ ». Chaque côté fonctionne comme une lunette
afocale dont l'image est redressée par des prismes.
\begin{questions}
\item Que signifient les nombres 10 et 50 ?
\item Le champ réel est l'angle du cône d'espace observé. Quelle largeur de paysage voit-on à \qty{1000}{m} ?
\item Le champ apparent est l'angle sous lequel l'œil voit cette image. Le calculer avec l'approximation des petits angles.
\item Le faisceau qui sort de chaque oculaire a un diamètre égal au diamètre de l'objectif divisé par le grossissement.
Le calculer et le comparer au diamètre de la pupille, 2 à \qty{7}{mm} selon l'éclairage.
\end{questions}
\end{exercice}
\end{blocLivre}
\rubrique{Objectif bac}
\begin{blocLivre}{ex-24-12}
\begin{exercice}[3]{La grande lunette de Meudon}
\begin{docu}{1}{Une lunette géante}
Inaugurée en 1891, la grande lunette de l'observatoire de Meudon possède un objectif de \qty{83}{cm} de diamètre et de
\qty{16,16}{m} de distance focale. L'astronome Eugène Antoniadi l'utilisa pour dessiner la surface de Mars.
\end{docu}
\begin{docu}{2}{Observer Mars}
Lors de ses oppositions les plus favorables, Mars apparaît sous un angle de $25''$ (secondes d'arc). La Lune est vue à
l'œil nu sous un angle de $0{,}52^\circ$.
\end{docu}
\begin{questions}
\item Calculer le grossissement obtenu avec un oculaire de \qty{25}{mm}.
\item Exprimer le diamètre apparent de Mars en radians, puis calculer l'angle sous lequel Mars apparaît à travers la
lunette, en degrés.
\item \initiative Quelle distance focale d'oculaire faudrait-il pour voir Mars aussi grande que la Lune à l'œil nu ?
\item Pourquoi les lunettes géantes ont-elles été abandonnées au profit des télescopes à miroir ?
\end{questions}
\end{exercice}
\end{blocLivre}
\begin{blocLivre}{ex-24-13}
\begin{exercice}[3]{Les satellites de Jupiter}
\begin{docu}{1}{La découverte de Galilée}
En janvier 1610, Galilée découvre avec sa lunette quatre satellites autour de Jupiter. Io, le plus proche, tourne à
\qty{4,22e5}{km} de la planète.
\end{docu}
\begin{docu}{2}{Les conditions d'observation}
À l'opposition, Jupiter est à \qty{6,3e8}{km} de la Terre et apparaît sous un diamètre de $45''$. L'éclat de Jupiter empêche
de voir Io à l'œil nu : il faut que l'écart apparent entre Io et le bord du disque de Jupiter dépasse environ $5'$
(minutes d'arc).
\end{docu}
\begin{questions}
\item Calculer l'écart angulaire maximal entre Io et le centre de Jupiter, en radians puis en minutes d'arc.
\item En déduire l'écart maximal entre Io et le bord du disque de Jupiter.
\item \initiative Quel grossissement minimal permet d'observer Io ? Des jumelles suffisent-elles ?
\item Galilée a vu les satellites changer de position d'une nuit à l'autre. Qu'en a-t-il déduit ?
\end{questions}
\end{exercice}
\end{blocLivre}
\begin{blocLivre}{ex-24-14}
\begin{exercice}[3]{Choisir ses oculaires}
\begin{docu}{1}{Une lunette commerciale}
Lunette 80/900 : objectif de \qty{80}{mm} de diamètre et de \qty{900}{mm} de distance focale. Livrée avec trois oculaires de
\qty{25}{mm}, \qty{12}{mm} et \qty{6}{mm}, et une lentille de Barlow « ×2 » qui double la distance focale de l'objectif.
\end{docu}
\begin{docu}{2}{Le conseil des astronomes}
« Le grossissement maximal utile vaut environ deux fois le diamètre de l'objectif en millimètres ; au-delà, l'image
s'agrandit mais devient floue et sombre. »
\end{docu}
\begin{questions}
\item Calculer les grossissements obtenus avec chaque oculaire, sans puis avec la lentille de Barlow.
\item Calculer le grossissement maximal utile de cette lunette.
\item \initiative Le fabricant annonce un grossissement de 300 sur l'emballage. Cette publicité est-elle honnête ?
\item Quelle combinaison conseiller pour observer les anneaux de Saturne, petits mais lumineux ?
\end{questions}
\end{exercice}
\end{blocLivre}
\begin{blocLivre}{ex-24-15}
\begin{exercice}[3]{La maquette de lunette}
\begin{docu}{1}{Le montage}
Sur un banc d'optique, on dispose d'un objectif de vergence \num{+2,0}~$\delta$ et d'un oculaire de vergence \num{+8,0}~$\delta$. Un
objet lumineux placé au foyer d'une lentille auxiliaire se comporte comme un objet à l'infini, vu sous un angle de
\qty{0,050}{rad}.
\end{docu}
\begin{docu}{2}{Les mesures}
Un écran placé à \qty{50}{cm} derrière l'objectif recueille une image nette de \qty{2,5}{cm}.
\end{docu}
\begin{questions}
\item Calculer les distances focales des deux lentilles et la distance qui doit les séparer.
\item Justifier la position de l'écran et la taille de l'image intermédiaire.
\item \initiative Retirer l'écran et calculer l'angle sous lequel l'œil voit l'image à travers l'oculaire. En déduire le
grossissement et le comparer à $f'_1/f'_2$.
\item Comment vérifier expérimentalement que le faisceau émergent est bien parallèle ?
\end{questions}
\end{exercice}
\end{blocLivre}
